\section{Literature Review and Expectations}\label{sec:literature}

The classical model treats evasion as a portfolio choice.
The taxpayer compares the tax saved with the penalty weighted by the probability of detection \citep{allingham1972}.
Later work widened that choice.
Evasion sits next to avoidance, and the audit rate is only one of the instruments the administration uses \citep{slemrod2007,slemrod2002}.
What the administration can see matters as much as the statutory penalty.
Where a third party reports the base, evasion is harder.
Where the taxpayer reports the base alone, it is easier \citep{kleven2011}.

The VAT makes that information problem concrete.
The invoice trail can deter evasion along the chain, because a buyer's credit is a seller's liability \citep{pomeranz2015}.
The trail does not close every gap.
Tax can be charged and then not remitted, credits can be misused, and an audit can still be the first time the administration records the amount \citep{keen2006}.
S\~ao Paulo has already been used to study one information instrument aimed at ICMS.
Consumer monitoring under the \emph{Nota Fiscal Paulista} raised reported activity \citep{naritomi2019}.
Work on the Brazilian state VAT has also studied fraudulent credits and the efficiency of aggregate collection \citep{gomes2024,cabello2020}.
Those papers stop at reporting, at fraud, or at the revenue the state books.
They do not describe the enrolled stock after the debt has been established, and they do not separate debts the firm declared from debts the administration assessed.

Beers, Hatch, Saldana and Wilson compare collection by the origin of the liability, mainly for individual income tax in the United States \citep{beers2015collectibility}.
The remaining question is how those origins are associated with the composition and the payment patterns of business VAT portfolios in other institutional settings.
A declared balance is a failure to pay a reported tax.
An assessed balance is a failure that the reporting system did not close, or a correction the administration imposed.
Both become certificates.
They need not have the same size, the same age, the same firms, or the same recovery.
A compliance study that treats enrolled \vat{} debt as one residual will average over those differences.
We do not test a model of why a firm ends in one door rather than the other.
We describe the two stocks so that later work can avoid pooling them by default.
Inside each segment, the partial associations of any payment and of the balance are the specification a linearity diagnostic selects.
Those associations do not explain how a firm arrives in one segment.

The description is organised by five expectations.
They are not hypotheses, and we do not treat a rejection as a test of the classical model.
First, if assessments are aimed at material understatements, assessed certificates should be larger than declared certificates.
Second, if declared debt is the unpaid remainder of routine filing, it should account for most certificates and a smaller share of value.
Third, a firm that has already reported the debt may be more likely to pay something, but the ratio of payments to the remaining stock need not rank the segments in the same way if assessed balances are much larger.
Fourth, a VAT stock dominated by large firms should be concentrated, so a small upper tail can hold most of the value.
Fifth, activity, age and registry status may differ across segments.
The firm register is observed after the debt, so those differences describe the debtors as they stand.
They do not identify a cause.
